
\section{Struktogramme und Klassendiagramme}

\textbf{Setze folgende Diagramme in BlueJ-Blockly um und notiere den Programmcode auf der Rückseite}

\vspace{0.5cm}

\textbf{Klassendiagramm:}

\includegraphics[width=0.5\textwidth]{Aufgaben/BlocksAndJava/img/klasse.png} 

\vspace{0.5cm}
\textbf{Struktogramm für piAusgeben(double x):}

\includegraphics[width=0.3\textwidth]{Aufgaben/BlocksAndJava/img/strukt_10pi.png}

\vspace{1cm}

\textbf{Zeichne hier die Struktogramme für die restlichen Methoden:}

\LoesungKaro{\\
\Large
+ double betrag(double z1) \\
\includegraphics[width=\textwidth]{Aufgaben/BlocksAndJava/img/strukt_betrag.png} \\\\
+ double potenzieren(double b, double p) \\
\includegraphics[width=\textwidth]{Aufgaben/BlocksAndJava/img/strukt_protenz.png}}{25}

\pagebreak

\Large
\begin{lstlisting}[escapechar=\&][]
public class ErweiterterTaschenrechner
{    
    public & \LoesungLuecke{double betrag(double z1) \{ }{15cm} &
    
        & \LoesungLuecke{if(z1 < 0) \{ }{15cm} &
        
            & \LoesungLuecke{z1 = -1 * z1;}{15cm} &
            
        & \LoesungLuecke{\}}{15cm} &        
    }
    public & \LoesungLuecke{double piAusgeben(double x) \{ }{15cm} &
    
        & \LoesungLuecke{for(int i=0; i < x; i++) \{ }{15cm} &
        
            & \LoesungLuecke{System.out.println(Math.PI);}{15cm} &
            
        & \LoesungLuecke{\}}{15cm} &
    }
    public & \LoesungLuecke{double potenzieren(double b, double p) \{ }{15cm} &

        & \LoesungLuecke{double ergebnis = 1.0;}{15cm} &
    
        & \LoesungLuecke{for(int i=0; i < p; i++) \{ }{15cm} &
        
            & \LoesungLuecke{ergebnis = ergebnis * b;}{15cm} &
            
        & \LoesungLuecke{\}}{15cm} &
            
        & \LoesungLuecke{return ergebnis;}{15cm} &
    }
}
\end{lstlisting}